Based on the computational synthesis and spectral plotting, the results highlight the classical trade-offs that complicate analog filter design. The comparative analysis of the four filter architectures against the predefined specifications suggests that no single approximation offers a flawless solution. Instead, each model forces a compromise between transition steepness, passband flatness, and phase linearity.

Examining the magnitude response reveals distinct behavioral priorities among the designs. The Butterworth filter maintained a maximally flat passband without introducing ripple artifacts. This smoothness, however, comes at a cost: the gradual roll-off dictates a considerably higher filter order to reach the required 60~dB of stopband attenuation, which may pose challenges in physical implementation due to component sensitivity. The Chebyshev approaches present an alternative strategy. They achieve a steeper transition by permitting either a 1~dB ripple in the passband (Type I) or equiripple variations in the stopband (Type II). Taking this mathematical compromise to the extreme, the Elliptic filter emerges as the most order-efficient architecture. By distributing ripples across both the passband and the stopband, it yields the sharpest magnitude transition of the four models, though this aggressive filtering introduces complications elsewhere.

When correlating these magnitude profiles with the corresponding phase analysis, the limitations of physical approximations become even more apparent. While the symbolic reference model maintains a theoretical zero-phase shift, every realizable filter introduced some degree of non-linear phase distortion. The Butterworth filter managed to sustain the most linear phase response within its passband. This characteristic is generally critical for minimizing temporal signal dispersion, particularly when precise wave shape preservation is required. On the other hand, models designed for sharp magnitude cutoffs, most notably the Elliptic and Chebyshev variations, exhibited severe phase non-linearities. These distortions appear most pronounced near the designated cutoff frequencies of 2000 and 4000~rad/sec, suggesting that applications requiring high phase fidelity might need to avoid these steeper filters entirely or implement additional phase equalization stages.

Ultimately, the experiment demonstrates that practical analog filter design is largely an exercise in managing competing parameters. Selecting a specific approximation cannot be reduced to a single metric of performance. Rather, it depends heavily on the primary demands of the target system. An engineer must weigh whether the application fundamentally requires absolute magnitude flatness, highly efficient transition steepness, or the strict phase linearity necessary to preserve signal timing and integrity.